% mct-worked-examples.tex — MCT 算例
% Source: examples/mct-worked-examples.md

\chapter{MCT 算例}\label{ex:mct}

对应章节：第 \ref{ch:mct} 章（MCT）。

\begin{examplebox}[title={例 1：RCT 正变换}]

\textbf{输入条件}：一个像素的 RGB 值，验证 RCT 正变换。

\medskip
\textbf{已知参数}：

\begin{tabular}{@{}lc@{}}
\toprule
参数 & 值 \\
\midrule
$R$ & 200 \\
$G$ & 128 \\
$B$ & 100 \\
位深 $d$ & 8 \\
$s_x, s_y$ & 1, 1 \\
\bottomrule
\end{tabular}

\medskip
\textbf{逐步计算}：

\begin{enumerate}
  \item $Y = (200 + 2 \times 128 + 100) \gg 2 = 556 \gg 2 = 139$
  \item $C_b = B - G = 100 - 128 = -28$
  \item $C_r = R - G = 200 - 128 = 72$
\end{enumerate}

\textbf{中间变量表}：

\begin{tabular}{@{}lc@{}}
\toprule
变量 & 值 \\
\midrule
$\mathit{tmp} = (R + 2G + B) \gg 2$ & 139 \\
$C_b = B - G$ & $-28$ \\
$C_r = R - G$ & 72 \\
$Y$ & 139 \\
\bottomrule
\end{tabular}

\medskip
\textbf{最终结果}：变换后分量值 $Y = 139, C_b = -28, C_r = 72$。分量 $c_0$（原 $R$）变为 139，$c_1$（原 $G$）变为 $-28$，$c_2$（原 $B$）变为 72。
\end{examplebox}

\begin{examplebox}[title={例 2：RCT 逆变换}]

\textbf{输入条件}：从例 1 的结果出发：$Y = 139, C_b = -28, C_r = 72$。

\medskip
\textbf{逐步计算}：

\begin{enumerate}
  \item $G = 139 - ((-28 + 72) \gg 2) = 139 - (44 \gg 2) = 139 - 11 = 128$
  \item $R = G + C_r = 128 + 72 = 200$
  \item $B = G + C_b = 128 + (-28) = 100$
\end{enumerate}

\medskip
\textbf{最终结果}：恢复后的值 $R = 200, G = 128, B = 100$。与输入完全一致。RCT 在整数域内是无损的。
\end{examplebox}

\begin{examplebox}[title={例 3：RCT 无损性证明}]

\textbf{输入条件}：$R = 1, G = 0, B = 0$，验证含舍入误差时的往返。

\medskip
\textbf{正变换}：
\begin{itemize}
  \item $Y = (1 + 0 + 0) \gg 2 = 1 \gg 2 = 0$
  \item $C_b = 0 - 0 = 0$
  \item $C_r = 1 - 0 = 1$
\end{itemize}

\textbf{逆变换}：
\begin{itemize}
  \item $G = 0 - ((0 + 1) \gg 2) = 0 - 0 = 0$
  \item $R = 0 + 1 = 1$
  \item $B = 0 + 0 = 0$
\end{itemize}

恢复正确。

\medskip
\textbf{证明}：设 $S = R + 2G + B$，则 $Y = \lfloor S/4 \rfloor$，$S \bmod 4$ 的信息可从 $C_b = B - G$ 和 $C_r = R - G$ 恢复：

\begin{align*}
S &= (G + C_r) + 2G + (G + C_b) = 4G + C_r + C_b \\
G &= Y - \lfloor (C_b + C_r) / 4 \rfloor \\
4G &= 4Y - 4\lfloor (C_b + C_r) / 4 \rfloor \\
S &= 4Y + (C_r + C_b) \bmod 4
\end{align*}

所以 RCT 是完全可逆的。
\end{examplebox}
